\documentclass[11pt,a4paper]{ctexart}
\usepackage[margin=1in]{geometry}
\usepackage{amsmath,amssymb}
\usepackage{graphicx}
\usepackage{booktabs}
\usepackage{array}
\usepackage{tabularx}
\usepackage{enumitem}
\usepackage{float}
\usepackage{hyperref}
\usepackage{xcolor}
\usepackage{caption}

\hypersetup{colorlinks=true,linkcolor=blue,urlcolor=blue,citecolor=blue}
\setlength{\parskip}{0.5em}
\setlength{\parindent}{2em}
\captionsetup{font=small}

\title{基于 MSWC English 精确首浊辅音子集的 /g/、/b/、/d/、/z/ 检测：理论、算法与案例分析\\\large VE216 Project Topic 2 重构版}
\date{2026年4月28日}

\begin{document}
\maketitle

\begin{abstract}
本文重构 VE216 Project Topic 2 的主线，不再把重点放在神经网络本身，而是放在任务定义、假设、理论分析、算法结构和定性案例上。我们把原先较松的词级多标签实验，收紧为一个更接近“首浊辅音检测”的四类任务：只保留首字母为 \texttt{g/b/d/z} 且 CMUdict 中首个非元音音素严格匹配 \texttt{G/B/D/Z} 的单词，共得到 60 个词、3000 条 16 kHz 语音。方法上，本文保留课内模板法作为主方法，并用 MLP 和 CNN 作为对照基线。结果表明，课程模板法可解释但误报高；MLP 的 exact-match 最好；CNN 的 macro-F1 最好。更重要的是，四类辅音的差异和声学事实一致：/z/ 最容易，/g/ 最难。
\end{abstract}

\section{问题定义与报告主线}

\subsection{任务形式化}

设输入语音为 $x_i$，目标标签为 $y_i \in \{0,1\}^4$，对应标签集合
\[
C=\{g,b,d,z\}.
\]
在当前数据集上，每个样本恰好属于四类中的一类，因此满足 $\sum_{c\in C} y_{i,c}=1$。但实现里仍保留一对多接口：先对每个标签输出分数 $s_{i,c}$，再用该标签自己的阈值 $\tau_c$ 做判定
\[
\hat y_{i,c} = \mathbb{1}[s_{i,c}\ge \tau_c].
\]
这样做的原因是：课内模板法、统计特征分类器和 CNN 输出的都不是天然校准过的概率，统一用 dev 集选阈值更公平，也更容易比较不同方法的特性。

\subsection{本文的重点}

这份报告只把 AI 模型当作对照基线，不把它们当成主角。主线是：
\begin{itemize}[nosep]
    \item 任务边界是否真的收紧到了“首浊辅音检测”；
    \item 课内模板法到底在假设什么、怎么计算、为什么会误报；
    \item 统计特征和时频特征为什么能比模板法更稳；
    \item /g/、/b/、/d/、/z/ 为什么会呈现出明显不同的难度。
\end{itemize}

\section{相关工作与理论背景}

本文关心的不是大词汇量识别，而是短语音里一个目标首辅音是否真的出现。这个问题更接近传统语音处理里的短时分析、局部谱形比较和模板匹配，而不是把整段音频直接交给一个黑盒分类器。

\subsection{短时分析与频谱表示}

Davis 和 Mermelstein 早就说明，语音识别里最有用的前端不是整段波形本身，而是经过短时分帧后的参数化表示\cite{davis1980}。做法很简单：把波形切成几十毫秒一帧，给每帧加窗，再做 Fourier transform，最后看这一小段时间里能量落在哪些频率上。这样得到的谱图，既保留了辅音释放、摩擦和元音过渡这些局部线索，也比原始波形稳定得多。

Rabiner 对 HMM 的综述则把这种帧级表示和序列建模连了起来\cite{rabiner1989}。它的核心意思是，语音不是静态图片，真正有意义的是“某个短时结构在时间上怎么移动”。虽然本文没有上 HMM，但这个观点很重要，因为它提醒我们：如果目标只是在整词里检出词首辅音，先把搜索范围压到短时块，再比较局部谱形，通常比直接看整词平均特征更合适。

\subsection{辅音线索为什么更依赖局部谱形}

对 /g/、/b/、/d/、/z/ 这类首辅音来说，最关键的证据往往只出现在很短的一段起始区间。爆破音的释放段短、转移快，摩擦音则在高频段拖得更久；这些差异一旦被整词平均掉，信息就会变少。Wilkinson 和 Russell 在 TIMIT 上做 phone recognition 时，也体现出同一个结论：越接近发音边界的局部线索，越能帮助辅音判别\cite{wilkinson2002}。

因此，本文的主方法没有先上复杂模型，而是回到课堂里最直接的工具链：短时 FFT、局部频谱、归一化相关。这样做的好处不是“更先进”，而是把问题本身讲清楚。

\subsection{本文为什么先用课堂方法}

本文后面的 AI 对照只负责补充结果，不负责定义主线。主线先从局部频谱模板出发，把“一个词首辅音在声学上长什么样”讲明白，再看更复杂的方法能把这个边界推到哪里。这样安排，后面的比较才不会变成单纯的分数表。


\section{数据集重建与标注规则}

\subsection{筛选规则}

新数据集由 \path{src/prepare_mswc_initial_gbdz_subset.py} 生成，核心规则只有一条：
\begin{quote}
保留一个单词，当且仅当它的首字母属于 \texttt{g/b/d/z}，并且 CMUdict / ARPABET 中首个非元音音素与该字母严格对应为 \texttt{G/B/D/Z}。
\end{quote}

这条规则的意义是把“字母里出现了某个浊辅音”改成“这个浊辅音确实出现在起始位置”。这样可以自动去掉两类噪声样本：
\begin{itemize}[nosep]
    \item 拼写里有目标字母，但发音位置不对；
    \item 首字母对了，但真正起始音素不是目标浊辅音。
\end{itemize}

\begin{verbatim}
for each word w in MSWC:
    if first letter of w not in {g,b,d,z}: discard
    phones = CMUdict(w)
    p = first non-vowel phone in phones
    if p == first letter(w).upper():
        keep w
    else:
        discard
\end{verbatim}

\subsection{数据规模}

最终实验规模如下：
\begin{table}[H]
    \centering
    \caption{精确首浊辅音子集的规模统计}
    \label{tab:dataset}
    \begin{tabular}{lcccc}
        \toprule
        标签 & 词数 & train & dev & test \\
        \midrule
        /g/ & 15 & 600 & 75 & 75 \\
        /b/ & 15 & 600 & 75 & 75 \\
        /d/ & 15 & 600 & 75 & 75 \\
        /z/ & 15 & 600 & 75 & 75 \\
        \midrule
        合计 & 60 & 2400 & 300 & 300 \\
        \bottomrule
    \end{tabular}
\end{table}

也就是说，这版数据是四类完全平衡的：总样本数 3000，测试集每类正样本 75 条。这个设定比旧版 600 词、多标签、\texttt{/b/\,/d/\,/l/\,/n/} 的版本更贴近题目，也更适合做案例分析。

\section{理论背景与方法}

\subsection{短时表示、局部窗与对比模板}

对 1 秒左右的短语音，真正有区分力的不是整词平均谱，而是词首浊辅音附近的局部时频结构。本文的前端仍然是标准短时分析：25 ms 帧长、10 ms 帧移、512 点 FFT、40 维 log-mel。不同的是，我们不再直接拿整词均值当模板，而是先把每帧 log-mel 压成能量曲线，平滑后用首次超过“基线 + margin”的位置当作 onset，再围绕 onset 截取固定长度窗口。

这个改造依赖两个假设：
\begin{itemize}[nosep]
    \item 目标辅音的判别证据主要集中在 onset 前后很短的一段时间里；
    \item 单独使用正模板不够，必须把“像目标”与“像干扰”同时建模。
\end{itemize}

因此，课内方法把每个样本变成两类局部描述：一类是窗口内分段均值拼接得到的 patch 向量，另一类是局部能量包络的自相关。训练时对每个标签分别学习正模板 $T_c^+$、负模板 $T_c^-$，以及正自相关模板 $A_c^+$、负自相关模板 $A_c^-$. 第 $i$ 个样本在平移 $k$ 下的得分定义为
\[
s_{i,c}^{(k)}=
0.65\bigl(\cos(p_i^{(k)},T_c^+)-\cos(p_i^{(k)},T_c^-)\bigr)
+
0.35\bigl(\cos(a_i^{(k)},A_c^+)-\cos(a_i^{(k)},A_c^-)\bigr).
\]
最终分数取候选平移上的最大值，$K=\{-2,0,2\}$。如果数据是单标签任务，就直接取 $\arg\max_c s_{i,c}$；如果保留一对多接口，则仍然可以回退到每类自己的阈值 $\tau_c$。

\subsection{算法 1：课程模板法（局部窗对比模板）}

\begin{verbatim}
Algorithm 1 局部窗对比模板法
Input: train/dev/test waveforms, labels C={g,b,d,z}
Output: per-label predictions on test

for each class c in C do
    collect positive and negative training samples
    build local patch templates T_c^+, T_c^- from onset-centered windows
    build local autocorr templates A_c^+, A_c^- from window envelopes
end for

for each sample i do
    detect onset from smoothed frame energy
    extract candidate windows around onset with shifts K={-2,0,2}
    for each class c in C do
        score each candidate by contrastive template similarity
        take the best candidate score as s_ic
    end for
    if single-label task then
        predict argmax_c s_ic
    else
        choose tau_c on dev and threshold
    end if
end for
\end{verbatim}

这个方法的核心不是“学一个复杂边界”，而是把判别问题改写为“局部窗里到底像不像目标辅音”。正样本模板提供典型形态，负样本模板提供干扰背景，两者之差比单独用正模板更稳。

\subsection{算法 2：统计特征 MLP 基线}

MLP 不是主角，只是一个强一点的统计基线。它的输入不是原始波形，而是手工汇总后的 62 维统计特征：
\begin{itemize}[nosep]
    \item 13 维 MFCC 的均值和标准差；
    \item MFCC 一阶差分的均值和标准差；
    \item 10 个 coarse band-energy 特征。
\end{itemize}

模型结构是 \texttt{StandardScaler -> MLPClassifier(128, 64)}，用 early stopping 训练。和模板法相比，它不再依赖显式模板对齐，而是直接学习这些统计特征之间的组合关系。因为当前数据是单标签四分类，但阈值仍按 label 单独选，所以 MLP 的 exact-match 往往会比 label-wise accuracy 更能反映它是否“少乱报”。

\subsection{算法 3：CNN 对照基线}

CNN 只作为更强的对照上界。输入是整词 log-mel spectrogram，网络是 3 个卷积块加自适应池化，再接一层小型全连接头。训练时使用带 \texttt{pos\_weight} 的 \texttt{BCEWithLogitsLoss}，并按 dev macro-F1 早停回退到最佳 epoch。

最终采用的调参配置是：
\[
\texttt{epochs}=10,\quad
\texttt{batch}=192,\quad
\texttt{lr}=10^{-3},\quad
\texttt{weight\_decay}=10^{-4},\quad
\texttt{patience}=5.
\]
它的作用只是提供一个“模型能学到多少局部时频模式”的参照，不影响本文的主要论证。

\section{定性案例分析}

\subsection{代表性样本}

图~\ref{fig:examples} 给出了四个代表词的波形和 spectrogram，图~\ref{fig:template} 给出了四类在训练集正样本上的平均 log-mel 模板。

\begin{figure}[H]
    \centering
    \includegraphics[width=0.95\textwidth]{figures/mswc_course_gbdz_local_patch_v1/letter_examples.png}
    \caption{四个目标浊辅音的代表性样本：\texttt{get}、\texttt{boy}、\texttt{did}、\texttt{zero}。}
    \label{fig:examples}
\end{figure}

\begin{figure}[H]
    \centering
    \includegraphics[width=0.90\textwidth]{figures/mswc_course_gbdz_local_patch_v1/template_grid.png}
    \caption{四个目标浊辅音在训练集正样本上的平均 log-mel 模板。}
    \label{fig:template}
\end{figure}

从图上能直接看出：/g/、/b/、/d/ 的关键差异主要集中在起始几十毫秒的短时爆破和过渡；/z/ 则有更长、更明显的高频摩擦能量。也因此，/z/ 的模板最稳定，/g/ 的模板最难。

\subsection{四个样本的预测行为}

\begin{table}[H]
    \centering
    \small
    \caption{四个代表词的预测行为与声学解释}
    \label{tab:cases}
    \begin{tabularx}{\textwidth}{l l l l l X}
        \toprule
        词 & 真值 & 课程模板法 & MLP & CNN & 解释 \\
        \midrule
        get & g & g & g & b & /g/ 的 burst 很短，局部窗抓住 onset 后 course 和 MLP 都能判对；CNN 仍容易和 /b/ 混淆。 \\
        boy & b & b & b & b & /b/ 的过渡最稳定，三种方法都比较容易。 \\
        did & d & z & d & b & /d/ 的 burst 更容易受对齐误差影响；course 方法在这个样本上仍会偏到 /z/，但 MLP 还能保住 stop 线索。 \\
        zero & z & z & z & d & /z/ 的摩擦噪声最长，course 和 MLP 都较稳；CNN 仍偶尔把它往 /d/ 拉。 \\
        \bottomrule
    \end{tabularx}
\end{table}

\begin{itemize}[nosep]
    \item 课程模板法经过局部窗和对比模板改造后，不再是整词均值判别；它的强项变成了 /b/、/d/、/z/ 这类 onset 线索清楚的样本。
    \item /g/ 仍最难，因为 burst 太短，稍微的对齐误差就会吞掉证据；这也是 MLP 在 /g/ 上略占优的原因。
    \item CNN 仍然最弱，说明在这个数据规模上，直接端到端看整词谱图并不能稳定学到首辅音证据。
\end{itemize}

\section{实验结果与比较}

\subsection{总体结果}

表~\ref{tab:summary} 是最终采用的 \texttt{local\_patch\_v1} 结果。局部窗改造后，课程模板法已经成为三种方法里最强的一个：它的 label-wise accuracy、exact-match 和 macro-F1 都最高，说明把判别证据限制在 onset 附近是有效的。MLP 仍然是稳定的统计基线，但整体略弱于课程模板法；CNN 在这一版设置下明显落后，说明仅靠整词级 log-mel 对短 burst 并不稳定。

\begin{table}[H]
    \centering
    \caption{总体指标比较。}
    \label{tab:summary}
    \begin{tabular}{lccccc}
        \toprule
        方法 & label-wise acc. & exact-match acc. & macro-F1 & micro-F1 & samples-F1 \\
        \midrule
        课程模板法 & \textbf{0.7750} & \textbf{0.5500} & \textbf{0.5462} & \textbf{0.5500} & \textbf{0.5500} \\
        MLP & 0.7633 & 0.5267 & 0.5258 & 0.5267 & 0.5267 \\
        CNN & 0.6600 & 0.3200 & 0.2219 & 0.3200 & 0.3200 \\
        \bottomrule
    \end{tabular}
\end{table}

\begin{figure}[H]
    \centering
    \includegraphics[width=0.92\textwidth]{figures/mswc_course_gbdz_local_patch_v1/metric_summary.png}
    \caption{三种方法的总体指标。}
    \label{fig:metric-summary}
\end{figure}

这里最重要的变化是：局部窗改造后，课程模板法的 exact-match 反而超过了 MLP；/b/、/d/、/z/ 的提升最明显，/g/ 仍然最难。CNN 的 macro-F1 最低，说明端到端看整词谱图并没有把词首证据稳定提出来。

\subsection{分标签结果}

\begin{table}[H]
    \centering
    \caption{四个标签的 F1}
    \label{tab:per-label}
    \begin{tabular}{lcccc}
        \toprule
        方法 & /g/ & /b/ & /d/ & /z/ \\
        \midrule
        课程模板法 & 0.438 & \textbf{0.528} & \textbf{0.472} & \textbf{0.747} \\
        MLP & \textbf{0.493} & 0.431 & 0.469 & 0.711 \\
        CNN & 0.049 & 0.444 & 0.024 & 0.371 \\
        \bottomrule
    \end{tabular}
\end{table}

\begin{figure}[H]
    \centering
    \includegraphics[width=0.90\textwidth]{figures/mswc_course_gbdz_local_patch_v1/per_label_f1.png}
    \caption{四个标签上的 F1 对比。}
    \label{fig:per-label-f1}
\end{figure}

从分标签看，/g/ 仍然是最难的；MLP 在 /g/ 上略占优，说明它学到了更多全局统计；课程模板法在 /b/、/d/、/z/ 上更稳，尤其是 /z/ 的持续摩擦噪声最容易被局部窗捕捉。CNN 在 /g/ 和 /d/ 上仍然偏弱，这和它更依赖大样本里稳定的局部纹理有关。

\section{讨论与局限}

这次重构最重要的收益，不是某个分数上升了，而是任务边界终于变清楚了。旧版更像“词里有没有出现某个浊辅音”，新版则更接近“词首这个浊辅音能不能被检出”。一旦把这个条件收紧，报告里的理论解释就顺了：数据集、特征、算法和案例分析都围绕同一个 onset 问题展开。

这套方法仍然有三个局限：
\begin{itemize}[nosep]
    \item 语音本身仍是整词级音频，不是人工切好的 phone segment；局部窗只是用 onset 近似对齐。
    \item /z/ 的词数仍然是数据集构造的瓶颈，想继续扩规模就得放宽单词池或换更大的语料。
    \item 目前的局部窗仍依赖阈值式 onset 检测，/g/ 和 /d/ 这类短爆破在对齐偏一点时就会掉分。
\end{itemize}

\section{结论}

本文完成了 Topic 2 的一次重新组织，核心结论如下：
\begin{itemize}[nosep]
    \item 通过“首字母与首个非元音音素严格对应”的规则，新数据集把任务收紧到了首浊辅音检测；
    \item 局部窗对比模板法在当前四类任务上取得最佳整体结果，macro-F1 为 0.546，exact-match 为 0.550；
    \item MLP 的 macro-F1 为 0.526，CNN 为 0.222，说明更强的深度模型并没有自动取代局部窗方法；
    \item /z/ 最容易、/g/ 最难，这和声学图里的持续摩擦与短促爆破完全一致；
    \item 这份报告的重点已经从“跑一个神经网络”转成“解释为什么这些辅音会这样表现”。
\end{itemize}

\begin{thebibliography}{99}
\bibitem{davis1980}
S. Davis and P. Mermelstein, ``Comparison of Parametric Representations for Monosyllabic Word Recognition in Continuously Spoken Sentences,'' \textit{IEEE Transactions on Acoustics, Speech, and Signal Processing}, 1980. \url{https://doi.org/10.1109/TASSP.1980.1163420}.

\bibitem{rabiner1989}
L. R. Rabiner, ``A Tutorial on Hidden Markov Models and Selected Applications in Speech Recognition,'' \textit{Proceedings of the IEEE}, 1989. \url{https://doi.org/10.1109/5.18626}.

\bibitem{wilkinson2002}
N. J. Wilkinson and M. J. Russell, ``Improved Phone Recognition on TIMIT Using Formant Frequency Data and Confidence Measures,'' in \textit{Proc. ICSLP 2002}. \url{https://www.isca-archive.org/icslp_2002/wilkinson02_icslp.html}.

\bibitem{chen2014kws}
G. Chen, C. Parada, and G. Heigold, ``Small-Footprint Keyword Spotting using Deep Neural Networks,'' in \textit{Proc. ICASSP 2014}. \url{https://research.google/pubs/small-footprint-keyword-spotting-using-deep-neural-networks/}.

\bibitem{sainath2015}
T. N. Sainath and C. Parada, ``Convolutional Neural Networks for Small-Footprint Keyword Spotting,'' in \textit{Proc. Interspeech 2015}. \url{https://research.google/pubs/convolutional-neural-networks-for-small-footprint-keyword-spotting/}.

\bibitem{arik2017}
S. O. Arik \textit{et al.}, ``Convolutional Recurrent Neural Networks for Small-Footprint Keyword Spotting,'' in \textit{Proc. Interspeech 2017}.

\bibitem{warden2018}
P. Warden, ``Speech Commands: A Dataset for Limited-Vocabulary Speech Recognition,'' arXiv:1804.03209, 2018. \url{https://doi.org/10.48550/arXiv.1804.03209}.

\bibitem{berg2021}
M. Berg \textit{et al.}, ``Keyword Transformer: A Self-Attention Model for Keyword Spotting,'' arXiv:2104.00769, 2021.

\bibitem{baevski2020}
A. Baevski \textit{et al.}, ``wav2vec 2.0: A Framework for Self-Supervised Learning of Speech Representations,'' arXiv:2006.11477, 2020. \url{https://doi.org/10.48550/arXiv.2006.11477}.

\bibitem{hsu2021hubert}
W.-N. Hsu \textit{et al.}, ``HuBERT: Self-Supervised Speech Representation Learning by Masked Prediction of Hidden Units,'' arXiv:2106.07447, 2021.

\bibitem{chen2022wavlm}
S. Chen \textit{et al.}, ``WavLM: Large-Scale Self-Supervised Pre-Training for Full Stack Speech Processing,'' \textit{IEEE Journal of Selected Topics in Signal Processing}, 2022. \url{https://arxiv.org/abs/2110.13900}.

\bibitem{gao2023}
C. Gao \textit{et al.}, ``Self-Supervised Speech Representation Learning for Keyword-Spotting with Light-Weight Transformers,'' arXiv:2303.04255, 2023. \url{https://arxiv.org/abs/2303.04255}.

\bibitem{kim2025}
D. Kim and J. Hwang, ``Fully End-to-end Streaming Open-vocabulary Keyword Spotting with W-CTC Forced Alignment,'' in \textit{Proc. Interspeech 2025}. \url{https://www.isca-archive.org/interspeech_2025/kim25d_interspeech.html}.
\end{thebibliography}

\end{document}
